\chapter{Aproximación por productos finitos de Blaschke}

Consideramos el espacio de Hardy $H^{\infty} \left(\mathbb{D}\right) \defeq \mathcal{H} \left(\mathbb{D}\right) \cap \mathcal{L}^\infty \left(\mathbb{D}\right)$. Por la \excref[prop:fn-holomorfa-esencialmente-acotada-imp-acotada]{prop-fn-holomorfa-esencialmente-acotada-imp-acotada}, tenemos que $H^{\infty} \left(\mathbb{D}\right) = \left\{f \in \mathcal{H} \left(\mathbb{D}\right) : f \tex{ está acotada en } \mathbb{D}\right\}$.

\begin{obs}\label{obs:hinfty-subespacio-vectorial-holomorfas}
    Como $\mathcal{H}\left(\mathbb{D}\right)$ es un espacio vectorial y $H^\infty$ es cerrado bajo suma y producto por escalares, $H^\infty \left(\mathbb{D}\right)$ es un subespacio vectorial de $\mathcal{H}\left(\mathbb{D}\right)$.
\end{obs}

\begin{prop}\label{prop:hinfty-esp-normado}
    La aplicación $\norm{\cdot}_{\infty} \colon H^\infty \left(\mathbb{D}\right) \to \mathbb{R}_{\geq 0}$ dada por $\norm{f}_{\infty} = \sup_{z \in \mathbb{D}} \abs{f(z)}$ es una norma en $H^\infty \left(\mathbb{D}\right)$. Además, $\forall f, g \in H^\infty \left(\mathbb{D}\right) : \norm{f \cdot g}_{\infty} \leq \norm{f}_{\infty} \norm{g}_{\infty}$.
\end{prop}
\begin{dem}
    Sean $f, g \in H^\infty \left(\mathbb{D}\right)$ y $\alpha \in \mathbb{C}$. Entonces,
    \begin{enumerate}[label=(\roman*)]
        \item $\ds \norm{f}_{\infty} = 0 \iff \sup_{z \in \mathbb{D}} \abs{f(z)} = 0 \iff \forall z \in \mathbb{D} : \abs{f(z)} = 0 \iff f = 0$.
        \item $\ds \norm{\alpha f}_{\infty} = \sup_{z \in \mathbb{D}} \abs{\alpha f(z)} = \abs{\alpha} \sup_{z \in \mathbb{D}} \abs{f(z)} = \abs{\alpha} \norm{f}_{\infty}$.
        \item $\begin{aligned}[t]
            \norm{f + g}_{\infty} &= \sup_{z \in \mathbb{D}} \abs{f(z) + g(z)} \leq \sup_{z \in \mathbb{D}} \left( \abs{f(z)} + \abs{g(z)} \right) \\
            &\leq \sup_{z \in \mathbb{D}} \abs{f(z)} + \sup_{z \in \mathbb{D}} \abs{g(z)} = \norm{f}_{\infty} + \norm{g}_{\infty}.
        \end{aligned}$
    \end{enumerate}

    Finalmente, $\ds \norm{f \cdot g}_{\infty} = \sup_{z \in \mathbb{D}} \abs{f(z) \cdot g(z)} \leq \left( \sup_{z \in \mathbb{D}} \abs{f(z)} \right) \left( \sup_{z \in \mathbb{D}} \abs{g(z)} \right) = \norm{f}_{\infty} \norm{g}_{\infty}$.
\end{dem}

\begin{cor}
    $H^\infty \left(\mathbb{D}\right)$ es un álgebra conmutativa con unidad.
\end{cor}
% \begin{dem}
%     Por la observación anterior, $H^\infty \left(\mathbb{D}\right)$ es un espacio vectorial.

%     Sean $f, g \in H^\infty \left(\mathbb{D}\right)$. Como $f$ y $g$ son holomorfas, $f \cdot g$ es holomorfa en $\mathbb{D}$. Además, por la \excref[prop:hinfty-esp-normado]{prop-hinfty-esp-normado}, $f \cdot g$ es acotada y, en consecuencia, $f \cdot g \in H^\infty \left(\mathbb{D}\right)$.

%     La conmutatividad, asociatividad y distributividad se heredan de las operaciones puntuales en $\mathbb{C}$. Finalmente, la función constante $1$ pertenece a $H^\infty \left(\mathbb{D}\right)$ y satisface $1 \cdot f = f$ para todo $f \in H^\infty \left(\mathbb{D}\right)$. \hfill\qedsymbol
% \end{dem}

\begin{obs}\label{obs:hinfty-bola-unidad}
    La bola unidad cerrada en $H^\infty \left(\mathbb{D}\right)$ según la métrica inducida por $\norm{\cdot}_{\infty}$ es precisamente la clase de Schur $\mathscr{S}$, es decir, $\left\{f \in H^\infty \left(\mathbb{D}\right) : \norm{f}_{\infty} \leq 1\right\} = \mathscr{S}$.
\end{obs}

\section{Aproximación de funciones en $\mathscr{S}$}

\begin{obs}
    Sean $f \in \mathscr{S}$ y $\left(B_n\right)_{n \in \N} \subset \exhyperref[defn:producto-finito-blaschke]{producto-finito-blaschke}{\mathscr{B}}$ tales que $B_n \xrightrightarrows[n \to \infty]{\mathbb{D}} f$. Como cada $B_n$ es holomorfa, está acotada por $1$ y es unimodular en $\mathbb{T}$, la convergencia uniforme asegura que $f$ se prolonga de forma continua y única a $\overline{\mathbb{D}}$ y que la restricción a $\mathbb{T}$ sigue siendo unimodular.
    
    Entonces, por el \excref[cor:carac-prod-finito-blaschke-algebra-disco]{cor-prod-finito-blaschke-algebra-disco}, $f \in \mathscr{B}$. En particular, $\mathscr{B}$ es un subconjunto de $\mathscr{S}$ cerrado respecto de la topología de la convergencia uniforme en $\mathbb{D}$. El siguiente ejemplo muestra que también se trata de un subconjunto propio.
\end{obs}

\begin{ejem}[de función en $\mathscr{S}$ sin extensión continua a $\bar{\mathbb{D}}$]
    Consideramos
    \[
    \forall z \in \mathbb{D} : f(z) = \exp \left(- \frac{1 + z}{1 - z} \right).
    \]
    Entonces, $f \in \mathcal{H} \left(\mathbb{D}\right)$ por ser composición de funciones holomorfas ($1 \notin \mathbb{D}$). Además,
    \[
    \forall z \in \mathbb{D} : \abs{f(z)} = \exp \left( - \Re \left( \frac{1 + z}{1 - z} \right) \right) \leq 1 \iff \Re \left( \frac{1 + z}{1 - z} \right) \geq 0.
    \]
    Sea $z = x + iy$ con $x^2 + y^2 < 1$. Entonces,
    \[
    \frac{1 + z}{1 - z} = \frac{(1+x+iy)(1-x+iy)}{(1-x)^2 + y^2} = \frac{(1-x^2-y^2) + 2iy}{(1-x)^2 + y^2}.
    \]
    Por lo tanto, $\Re \left( \frac{1 + z}{1 - z} \right) = \frac{1 - x^2 - y^2}{(1-x)^2 + y^2} > 0$, ya que el numerador es positivo (pues $x^2 + y^2 < 1$) y el denominador es positivo. Así, $\forall z \in \mathbb{D} : \abs{f(z)} < 1$, lo que demuestra que $f \in \mathscr{S}$.

    Sin embargo, $f$ no se puede prolongar de forma continua a $\bar{\mathbb{D}}$. Para verlo, consideramos dos caminos que tienden a $1$:
    \begin{itemize}
        \item $\ds \lim_{t \to 1^-} f(t) = \lim_{t \to 1^-} \exp \left( - \frac{1 + t}{1 - t} \right) = \exp(-\infty) = 0$.
        \item $\ds \lim_{t \to 0} f \left(e^{it} \right) = \lim_{t \to 0} \exp \left( - \frac{1 + e^{it}}{1 - e^{it}} \right) = \lim_{t \to 0} \exp \left( - i \cot \left( \frac{t}{2} \right) \right) = 1$.
    \end{itemize}
    Como los límites difieren, $f$ no se puede prolongar de forma continua a $\bar{\mathbb{D}}$.
\end{ejem}

Sin embargo, si nos limitamos a subconjuntos compactos de $\mathbb{D}$, podemos aproximar cualquier función en $\mathscr{S}$ por productos finitos de Blaschke. En otras palabras, $\exhyperref[defn:producto-finito-blaschke]{producto-finito-blaschke}{\mathscr{B}}$ es denso en $\mathscr{S}$ para la topología de la convergencia uniforme sobre compactos de $\mathbb{D}$.

\begin{lem}\label{lem:aprox-clase-schur-prod-finito-blaschke-orden-cero}
    $\forall n \in \N, f \in \exhyperref[defn:clase-schur]{clase-schur}{\mathscr{S}} : \exists B_n \in \exhyperref[defn:producto-finito-blaschke]{producto-finito-blaschke}{\mathscr{B}} : \exhyperref[defn:orden-cero-fn-holomorfa]{orden-cero-fn-holomorfa}{\operatorname{ord}}(f - B_n, 0) \geq n$.
\end{lem}
\begin{dem}
    Dados $n \in \N$ y $f \in \mathscr{S}$, si $\abs{f(0)} = 1$, por el \exhyperref[cor:modulo-maximo]{cor-modulo-maximo}{principio del módulo máximo}, $f$ es una constante unimodular y, por tanto, $f \in \mathscr{B}$. En consecuencia, basta tomar $B_n = f$. Supongamos entonces que $f(0) \in \mathbb{D}$.

    Lo probaremos por inducción en $n$. Para $n = 1$ y $f \in \mathscr{S}$ arbitraria, basta tomar $B_1 = \exhyperref[defn:involucion-disco-unidad]{involucion-disco-unidad}{\tau_{f(0)}} \in \mathscr{B}$, ya que $(f - B_1)(0) = f(0) - \tau_{f(0)}(0) = f(0) - f(0) = 0$.

    Sea $n \in \N$, supongamos que $\forall f \in \mathscr{S} : \exists B_n \in \mathscr{B} : \operatorname{ord}(f - B_n, 0) \geq n$. Sea $f \in \mathscr{S}$ y escribimos $c_0 \defeq f(0)$. Consideramos $g \colon \mathbb{D} \to \C$ dada por
    \begin{equation}\label{dem:aprox-clase-schur-prod-finito-blaschke-orden-cero:ecu:g}
        \forall z \in \mathbb{D} \setminus \{0\} : g(z) \defeq \frac{\tau_{c_0}(f(z))}{z}.
    \end{equation}
    Por un lado, como $\tau_{c_0} (f(0)) = 0$, por el \excref[teo:singularidad-evitable-riemann]{teo-singularidad-evitable-riemann}, $g$ se puede prolongar de forma holomorfa a todo $\mathbb{D}$. Por otro lado, como $f \in \mathscr{S}$ y $\tau_{c_0} \in \mathrm{Aut} (\mathbb{D})$, se tiene que $\tau_{c_0} \circ f \in \mathscr{S}$ y el \exhyperref[lem:schwarz]{lem-schwarz}{lema de Schwarz} nos garantiza que $\forall z \in \mathbb{D} : \abs{\tau_{c_0} (f(z))} \leq \abs{z}$. En consecuencia, $\forall z \in \mathbb{D} : \abs{g(z)} \leq 1$, es decir, $g \in \mathscr{S}$.

    Por la hipótesis de inducción, $\exists B_n \in \mathscr{B} : \operatorname{ord}(g - B_n, 0) \geq n$. Es decir, $\exists G_1 \in \mathcal{H}(\mathbb{D})$ tal que $\forall z \in \mathbb{D} : g(z) - B_n(z) = z^n G_1(z)$. Además, como $g, B_n \in \mathscr{S}$, se tiene que $\norm{g - B_n}_{\infty} \leq 2$. Fijamos $r = \sfrac{1}{2}$. Entonces:
    \begin{itemize}
        \item Si $\abs{z} \geq r$, $\ds \abs{G_1(z)} = \frac{\abs{g(z) - B_n(z)}}{\abs{z}^n} \leq \frac{2}{r^n}$.
        \item Si $\abs{z} < r$, por el \exhyperref[cor:modulo-maximo]{cor-modulo-maximo}{principio del módulo máximo}, $\ds \abs{G_1(z)} \leq \max_{\abs{w}=r} \abs{G_1(w)} \leq \frac{2}{r^n}$.
    \end{itemize}
    Por tanto, $G_1$ es acotada en $\mathbb{D}$ y $G_1 \in H^\infty\left(\mathbb{D}\right)$.

    Definimos $\forall z \in \mathbb{D} : B_{n+1}(z) \defeq \tau_{c_0}\left( z B_n(z) \right)$. Entonces, $B_{n+1} \in \mathscr{B}$ por el \excref[lem:prod-finito-blaschke-composicion-involucion]{lem-prod-finito-blaschke-composicion-involucion}. De \eqref{dem:aprox-clase-schur-prod-finito-blaschke-orden-cero:ecu:g}, tenemos que $f(z) = \tau_{c_0}^{-1} (z g(z)) = \tau_{c_0} (z g(z))$ aplicando el \excref[lem:involucion-disco-unidad]{lem-involucion-disco-unidad}. Por tanto, obtenemos que $\forall z \in \mathbb{D} : $
    \[
    \begin{aligned}
        f(z) - B_{n+1}(z) &= \tau_{c_0} (z g(z)) - \tau_{c_0} (z B_n(z)) = \frac{c_0 - z g(z)}{1 - \bar{c_0} z g(z)} - \frac{c_0 - z B_n(z)}{1 - \bar{c_0} z B_n(z)}
        \\
        % &= \frac{(c_0 - z g(z))(1 - \bar{c_0} z B_n(z)) - (c_0 - z B_n(z))(1 - \bar{c_0} z g(z))}{(1 - \bar{c_0} z g(z)) (1 - \bar{c_0} z B_n(z))} \\
        % &= \frac{c_0 - z g(z) - \bar{c_0} c_0 z B_n(z) + \bar{c_0} z^2 g(z) B_n(z) - c_0 + z B_n(z) + \bar{c_0} c_0 z g(z) - \bar{c_0} z^2 g(z) B_n(z)}{(1 - \bar{c_0} z g(z)) (1 - \bar{c_0} z B_n(z))} \\
        % &= \frac{z (g(z) - B_n(z)) + \bar{c_0} c_0 z (g(z) - B_n(z))}{(1 - \bar{c_0} z g(z)) (1 - \bar{c_0} z B_n(z))} \\
        &= \underbrace{\frac{\left(1 - \abs{c_0}^2 \right)}{\left(1 - \bar{c_0} z g(z) \right) \left(1 - \bar{c_0} z B_n(z) \right)}}_{\eqdef G_2(z)} \cdot z \left( g(z) - B_n(z) \right) = G_2(z)\, z^{n+1} G_1(z).
    \end{aligned}
    \]
    Además, como $\abs{g(z)} \leq 1$ y $\abs{B_n(z)} \leq 1$, por la desigualdad triangular inversa 
    \[
        \forall z \in \mathbb{D} : \abs{1 - \bar{c_0} z g(z)} \geq 1 - \abs{c_0}\abs{z}\abs{g(z)} > 1 - \abs{c_0} > 0,
    \]
    y análogamente $\abs{1 - \bar{c_0} z B_n(z)} > 0$. Luego $G_2 \in \mathcal{H}(\mathbb{D})$, y por tanto $G_2G_1 \in \mathcal{H}(\mathbb{D})$. En consecuencia, $\operatorname{ord}(f - B_{n+1}, 0) \geq n+1$.\hfill\qedsymbol
\end{dem}

\begin{teo}[Carathéodory]\label{teo:caratheodory-aproximacion-schur-prod-finito-blaschke-compactos}
    $\forall f \in \exhyperref[defn:clase-schur]{clase-schur}{\mathscr{S}} : \exists \left(B_n\right)_{n \in \mathbb{N}} \subset \exhyperref[defn:producto-finito-blaschke]{producto-finito-blaschke}{\mathscr{B}} : B_n \xRightarrow[n \to \infty]{\mathbb{D}} f$. 
\end{teo}
\begin{dem}
    Sea $f \in \mathscr{S}$. Por el \excref[lem:aprox-clase-schur-prod-finito-blaschke-orden-cero]{lem-aprox-clase-schur-prod-finito-blaschke-orden-cero}, $\forall n \in \N : \exists B_n \in \mathscr{B} : \operatorname{ord}(f - B_n, 0) \geq n$. Es decir, dado $n \in \N$, $\exists g_n \in H^\infty : \forall z \in \mathbb{D} : f(z) - B_n(z) = z^{n-1} g_n(z) \we g_n(0) = 0$.
    
    Además, $\norm{z^{n-1}g_n}_{\infty} \leq \norm{g_n}_{\infty}$, ya que $\forall z \in \mathbb{D} : \abs{z}^{n-1} \leq 1$. Veamos la desigualdad recíproca. Tomamos una sucesión maximizante $\left(z_m\right)_{m \in \N} \subset \mathbb{D}$ de $g_n$, es decir,
    \[
    \abs{g_n(z_m)} \xrightarrow[m \to \infty]{} \norm{g_n}_{\infty}.
    \]
    Supongamos $\norm{g_n}_{\infty} > 0$ (si $\norm{g_n}_{\infty} = 0$, no hay nada que probar). Sea $\left(z_{m_j}\right)_{j \in \N}$ una subsucesión convergente a $\zeta \in \overline{\mathbb{D}}$. Si $\abs{\zeta} < 1$, por continuidad se tendría
    \[
    \abs{g_n(\zeta)} = \lim_{j \to \infty} \abs{g_n(z_{m_j})} = \norm{g_n}_{\infty},
    \]
    luego, por el \exhyperref[teo:modulo-maximo]{teo-modulo-maximo}{principio del módulo máximo}, $g_n$ sería constante en $\mathbb{D}$; y como $g_n(0) = 0$, esto contradice que $\norm{g_n}_{\infty} > 0$. Por tanto, $\abs{\zeta} = 1$ y $\smash{\abs{z_{m_j}}^{n-1} \xrightarrow[j \to \infty]{} 1}$. Por lo tanto,
    \[
    \norm{g_n}_{\infty} = \lim_{j \to \infty} \abs{g_n(z_{m_j})} = \lim_{j \to \infty} \frac{\abs{z_{m_j}^{n-1} g_n(z_{m_j})}}{\abs{z_{m_j}}^{n-1}} \leq \norm{z^{n-1} g_n}_{\infty}.
    \]
    Así, $\norm{g_n}_{\infty} = \norm{z^{n-1} g_n}_{\infty}$. Por otro lado, como $f, B_n \in \mathscr{S}$, $\norm{f - B_n}_{\infty} \leq 2$. Por lo tanto, $\norm{g_n}_{\infty} = \norm{z^{n-1} g_n}_{\infty} = \norm{f - B_n}_{\infty} \leq 2$.

    Entonces, $h_n \defeq \sfrac{g_n}{2} \in \mathscr{S}$ y $h_n(0) = 0$, luego, por el \exhyperref[lem:schwarz]{lem-schwarz}{lema de Schwarz} aplicado a $h_n$, tenemos que $\forall z \in \mathbb{D} : \abs{h_n(z)} \leq \abs{z}$. En consecuencia,
    \[
    \forall z \in \mathbb{D} : \abs{f(z) - B_n (z)} = \abs{z^{n-1} g_n(z)} = 2 \abs{z^{n-1} h_n(z)} \leq 2 \abs{z}^n.
    \]
    Por tanto, para $K \subset \mathbb{D}$ compacto, existe $r \in (0, 1)$ tal que $\forall z \in K : \abs{z} \leq r$. Así, $\sup_{z \in K} \abs{f(z) - B_n(z)} \leq 2 r^n \xrightarrow[n \to \infty]{} 0$, i.e. $B_n \xrightrightarrows[n \to \infty]{K} f$.\hfill\qedsymbol
\end{dem}

\section{La envolvente convexa de $\mathscr{B}$}

\begin{lem}\label{lem:prod-combinacion-convexa-prod-finito-blaschke}
    Sean $f, g \in \exhyperref[defn:envolvente-convexa]{envolvente-convexa}{\operatorname{conv}} \left(\exhyperref[defn:producto-finito-blaschke]{producto-finito-blaschke}{\mathscr{B}}\right)$. Entonces, $f \cdot g \in \exhyperref[defn:envolvente-convexa]{envolvente-convexa}{\operatorname{conv}} \left(\exhyperref[defn:producto-finito-blaschke]{producto-finito-blaschke}{\mathscr{B}}\right)$.
\end{lem}
\begin{dem}
    Por la \excref[prop:envolvente-convexa-combinacion-convexa]{prop-envolvente-convexa-combinacion-convexa}, sabemos que $\exists \left\{f_k\right\}_{k=1}^n, \left\{g_\ell\right\}_{\ell=1}^m \subset \mathscr{B}$ y $\left\{\alpha_k\right\}_{k=1}^n, \left\{\beta_\ell\right\}_{\ell=1}^m \subset \R_{\geq 0} : \sum_{k} \alpha_k = 1 \we \sum_{\ell} \beta_\ell = 1 \we f = \sum_{k} \alpha_k f_k \we g = \sum_{\ell} \beta_\ell g_\ell$.
    \[
    \therefore \quad f \cdot g = \left( \sum_{k=1}^n \alpha_k f_k \right) \cdot \left( \sum_{\ell=1}^m \beta_\ell g_\ell \right) = \sum_{k=1}^n \sum_{\ell=1}^m \alpha_k \beta_\ell \, f_k \cdot g_\ell.
    \]
    Como $f_k, g_\ell \in \mathscr{B}$, $f_k \cdot g_\ell \in \mathscr{B}$. Además, $\sum_{k} \sum_{\ell} \alpha_k \beta_\ell = \left( \sum_{k} \alpha_k \right) \left( \sum_{\ell} \beta_\ell \right) = 1$. Por lo tanto, $f \cdot g$ es una combinación convexa de productos finitos de Blaschke.\hfill\qedsymbol
\end{dem}%ORDENAR enviar al anexo

\begin{lem}\label{lem:dilatacion-prod-finito-blaschke}
    Sean $B \in \exhyperref[defn:producto-finito-blaschke]{producto-finito-blaschke}{\mathscr{B}}$ y $t \in [0, 1)$. Entonces, $B_t$ dado por $\forall z \in \mathbb{D} : B_t(z) = B(tz)$ es una \exhyperref[defn:combinacion-convexa]{combinacion-convexa}{combinación convexa} de \exhyperref[defn:producto-finito-blaschke]{producto-finito-blaschke}{productos finitos de Blaschke}.
\end{lem}
\begin{dem}
    Sea $B = \gamma \prod_{k=1}^n \exhyperref[defn:involucion-disco-unidad]{involucion-disco-unidad}{\tau_{w_k}}$ con $\gamma \in \mathbb{T}$, $n \in \N$ y $w_1, \dots, w_n \in \mathbb{D}$. Entonces, $B_t = \gamma \prod_{k=1}^n (\tau_{w_k})_t$ y, como por el \excref[lem:prod-combinacion-convexa-prod-finito-blaschke]{lem-prod-combinacion-convexa-prod-finito-blaschke}, el producto de combinaciones convexas de productos finitos de Blaschke es una combinación convexa de productos finitos de Blaschke, basta probar que el lema se cumple para $B = \tau_w$ con $w \in \mathbb{D}$.

    Si $w = 0$, el resultado es inmediato. Suponemos $w \neq 0$. Escribimos $\tau_w(tz)$ como una combinación lineal de $\tau_{wt}(z)$ y una constante:
    \[
    c_1 \tau_{wt}(z) + \kappa = c_1 \frac{wt - z}{1 - \overline{w} t z} + \kappa = (\tau_w)_t(z) = \frac{w - t z}{1 - \overline{w} t z}.
    \]
    Resolviendo el sistema de ecuaciones resultante, obtenemos que
    \[
    % \begin{cases}
    %     c_1 wt + \kappa = w \\
    %     \bar{w}t \kappa + c_1 = t
    % \end{cases}
    % \implies
    c_1 = \frac{t(1 - \abs{w}^2)}{1 - \abs{w}^2 t^2} \quad \we \quad \kappa = \frac{w(1 - t^2)}{1 - \abs{w}^2 t^2} = \frac{\abs{w}(1 - t^2)}{1 - \abs{w}^2 t^2} e^{i \arg w}.
    \]
    Definiendo $c_2 \defeq \frac{\abs{w}(1 - t^2)}{1 - \abs{w}^2 t^2}$, $(\tau_w)_t$ queda escrito como combinación lineal de dos productos finitos de Blaschke: $\tau_{wt}, \exhyperref[ejem:rotacion-disco-unidad]{ejem-rotacion-disco-unidad}{\rho_{\arg w}} \in \mathscr{B}$.
    Ahora bien, esta no es una combinación convexa, pues la suma de los coeficientes es
    \[
    c_1 + c_2 = \frac{t - t\abs{w}^2 + \abs{w} - \abs{w}t^2}{(1 - \abs{w}t)(1 + \abs{w}t)} = \frac{(t + \abs{w})\cancel{(1 - t\abs{w})}}{\cancel{(1 - \abs{w}t)}(1 + \abs{w}t)} = \frac{t + \abs{w}}{1 + \abs{w}t} < 1.
    \]
    El defecto es $\delta = 1 - (c_1 + c_2) = \frac{(1-t)(1-\abs{w})}{1 + \abs{w}t}$, por lo que si añadimos los términos $\frac{\delta}{2} \cdot 1 + \frac{\delta}{2} \cdot (-1)$, completamos la combinación convexa:
    \[
    \begin{aligned}
        {(\tau_w)}_t &= c_1 \tau_{wt} + c_2 \rho_{\arg w} + \frac{\delta}{2} \rho_{0} + \frac{\delta}{2} \rho_{\pi} \\
        &= \frac{t(1 - \abs{w}^2)}{1 - \abs{w}^2 t^2} \tau_{wt} + \frac{\abs{w}(1 - t^2)}{1 - \abs{w}^2 t^2} \rho_{\arg w} + \frac{(1-t)(1-\abs{w})}{2(1 + \abs{w}t)} \rho_{0} + \frac{(1-t)(1-\abs{w})}{2(1 + \abs{w}t)} \rho_{\pi},
    \end{aligned}
    \]
    donde $c_1, c_2, \frac{\delta}{2} \geq 0$ y $c_1 + c_2 + \frac{\delta}{2} + \frac{\delta}{2} = 1$. \hfill\qedsymbol
\end{dem}

\begin{teo}[Fisher]\label{teo:fisher-aproximacion-algebra-disco-conv-prod-finito-blaschke}
    Sea $f \in \exhyperref[defn:algebra-disco-unidad]{algebra-disco-unidad}{\mathscr{A}} \left(\mathbb{D} \right)$ tal que $\exhyperref[defn:norma-lp]{norma-lp}{\norm{f}_{\infty}} \leq 1$. Entonces,
    \[
    \exists \left(f_n\right)_{n\in \N} \subset \exhyperref[defn:envolvente-convexa]{envolvente-convexa}{\operatorname{conv}} (\exhyperref[defn:producto-finito-blaschke]{producto-finito-blaschke}{\mathscr{B}}) : f_n \exhyperref[defn:convergencia-uniforme]{convergencia-uniforme}{\xrightrightarrows[n \to \infty]{\overline{\mathbb{D}}}} f.
    \]
\end{teo}
\begin{dem}
    Sea $\varepsilon > 0$. Como $f \in \mathcal{C} \left(\bar{\mathbb{D}}\right)$ y $\bar{\mathbb{D}}$ es compacto, $f$ es uniformemente continua en $\overline{\mathbb{D}}$. Por lo tanto, $\exists t \in [0, 1) : \max_{z \in \overline{\mathbb{D}}} \abs{f(tz) - f(z)} < \sfrac{\varepsilon}{2}$.

    Como $\overline{t\mathbb{D}} \Subset \mathbb{D}$ y $f \in \mathscr{S}$, por el \excref[teo:caratheodory-aproximacion-schur-prod-finito-blaschke-compactos]{teo-caratheodory-aproximacion-schur-prod-finito-blaschke-compactos}, $\exists B \in \mathscr{B} :\ds \smash{\max_{w \in \overline{t\mathbb{D}}} \abs{f(w) - B(w)} < \sfrac{\varepsilon}{2}}$. Definiendo $f_t(z) \defeq f(tz)$ y $B_t(z) \defeq B(tz)$,
    \[
    \max_{z \in \overline{\mathbb{D}}} \abs{f_t(z) - B_t(z)} = \max_{z \in \overline{\mathbb{D}}} \abs{f(tz) - B(tz)} = \max_{w \in \overline{t\mathbb{D}}} \abs{f(w) - B(w)} < \frac{\varepsilon}{2}.
    \]
    En consecuencia, $\ds \max_{z \in \overline{\mathbb{D}}} \abs{f(z) - B_t(z)} \leq \max_{z \in \overline{\mathbb{D}}} \abs{f(z) - f_t(z)} + \max_{z \in \overline{\mathbb{D}}} \abs{f_t(z) - B_t(z)} < \varepsilon$.

    Además, por el \excref[lem:dilatacion-prod-finito-blaschke]{lem-dilatacion-prod-finito-blaschke}, $B_t$ es una combinación convexa de productos finitos de Blaschke. Como $\varepsilon > 0$ es arbitrario, tomando $\varepsilon = \sfrac{1}{n}$ para $n \in \N$ obtenemos una sucesión $\left(f_n\right)_{n \in \N} \subset \operatorname{conv}(\mathscr{B})$ tal que $f_n$ converge uniformemente a $f$ en $\overline{\mathbb{D}}$.\hfill\qedsymbol
\end{dem}

\section{Aproximación de funciones unimodulares continuas}

Recordemos que para $f$ meromorfa en $\mathbb{D}$ tal que $f \in \mathcal{C}\left(\bar{\mathbb{D}}\right)$ y $f (\mathbb{T}) \subset \mathbb{T}$, el \excref[cor:carac-cociente-prod-finito-blaschke-algebra-disco-meromorfa]{cor-carac-cociente-prod-finito-blaschke-algebra-disco-meromorfa} nos garantiza que existen $B_1, B_2 \in \mathscr{B}$ tal que $f = \sfrac{B_1}{B_2}$.

Si nos restringimos a la frontera $\mathbb{T}$, podemos rebajar las hipótesis y demostrar un resultado de aproximación. Para ello, necesitamos el siguiente lema, que esencialmente afirma que una curva cerrada que no pasa por el origen da un número par o impar de vueltas alrededor de este.

\begin{lem}\label{lem:fn-continua-toro-imp-vueltas-par-impar}
    Sea $\gamma \colon \mathbb{T} \to \mathbb{T}$ \exhyperref[defn:continuidad-pnt]{continuidad}{continua}. Entonces, $\exists \eta \in \mathcal{C}\left(\mathbb{T} \right)$ unimodular tal que
    \[
    \left(\forall \zeta \in \mathbb{T} : \gamma(\zeta) = \eta^2(\zeta) \right) \vebar \left(\forall \zeta \in \mathbb{T} : \gamma(\zeta) = \zeta \cdot \eta^2(\zeta) \right).
    \] 
\end{lem}
\begin{dem}%¿No se podría demostrar por argumentos topológicos?
    Como $\gamma \in \mathcal{C}\left(\mathbb{T} \right)$ y $\mathbb{T}$ es compacto, $\gamma$ es uniformemente continua en $\mathbb{T}$. Es decir, $\exists N \in \N : \forall \theta, \theta' \in \R : \abs{\theta - \theta'} \leq \frac{2\pi}{N} \implies \abs{\gamma \left(e^{i\theta} \right) - \gamma \left(e^{i\theta'} \right)} < 2$.
    
    Dividimos $[0, 2\pi]$ y, por tanto, la circunferencia $\mathbb{T}$ en $N$ arcos de longitud $\frac{2\pi}{N}$:
    \[
    \begin{aligned}
        [0, 2\pi] &= \bigcup_{k=1}^N I_k \quad \tex{donde}\quad \forall k \in \N_N : I_k = \left[\frac{2\pi (k-1)}{N}, \frac{2\pi k}{N} \right] \\
        \implies \mathbb{T} &= \bigcup_{k=1}^N T_k \quad \tex{donde}\quad \forall k \in \N_N : T_k = \left\{ e^{i\theta} : \theta \in I_k \right\}.
    \end{aligned}
    \]
    Como $\gamma$ es continua y $\forall k \in \N_N : \forall \zeta, \zeta' \in T_k : \abs{\gamma(\zeta) - \gamma(\zeta')} < 2$ y $T_k$ es compacto y conexo, $\gamma(T_k)$ es un arco cerrado de $\mathbb{T}$ que abarca menos de $\pi$ radianes de ángulo.

    Por tanto, $\forall k \in \N_N : \exists \phi_k \colon I_k \to \R$ continua tal que $\forall \theta \in I_k : \gamma\left(\theta \right) = e^{i\phi_k(\theta)}$. Cada $\phi_k$ está unívocamente determinada salvo por la suma de un múltiplo entero de $2\pi$. Como $\gamma$ es continua, podemos elegir las funciones $\phi_k$ de forma que $\phi_k \left( \frac{2k\pi}{N} \right) = \phi_{k+1} \left( \frac{2k\pi}{N} \right)$ para todo $k \in \N_{N-1}$. Así, la aplicación $\phi \colon \R \to \R$ dada por $\forall k \in \N_N : \forall \theta \in I_k : \phi(\theta) = \phi_k(\theta)$ es continua y, además, por la continuidad de $\gamma$, $\exists m \in \Z : \phi(2\pi) = \phi(0) + 2m\pi$.

    Entonces, se tiene uno de estos dos casos:
    \begin{enumerate}[label=\Roman*.]
        \item Si $m$ es par, definimos $\forall \zeta = e^{i\theta} \in \mathbb{T} : \eta(\zeta) = \exp \left( i \frac{\phi(\theta)}{2} \right)$. Entonces, $\eta \in \mathcal{C}\left(\mathbb{T} \right)$ y $\forall \zeta = e^{i\theta} \in \mathbb{T} : \eta^2(\zeta) = e^{i\phi(\theta)} = \gamma(\zeta)$.
        \item Si $m$ es impar, definimos $\forall \zeta = e^{i\theta} \in \mathbb{T} : \eta(\zeta) = \exp \left( i \frac{1}{2} \left(\phi(\theta) - \theta\right)\right)$. Entonces, $\eta \in \mathcal{C}\left(\mathbb{T} \right)$ y $\forall \zeta = e^{i\theta} \in \mathbb{T} : \zeta \cdot \eta^2(\zeta) = e^{i\theta} \cdot \exp \left( i \left(\phi(\theta) - \theta\right) \right) = e^{i\phi(\theta)} = \gamma(\zeta)$.
    \end{enumerate}

    \vspace{-\baselineskip}
    % \vspace{-\belowdisplayskip}
    \hfill\qedsymbol
\end{dem}

\begin{teo}[Helson--Sarason]\label{teo:helson-sarason-aproximacion-fn-unimodular-continua-cociente-prod-finito-blaschke}
    Sea $u \colon \mathbb{T} \to \mathbb{T}$ \exhyperref[defn:continuidad-pnt]{continuidad}{continua} y sea $\varepsilon > 0$. Entonces, existen $B_1, B_2$ \exhyperref[defn:producto-finito-blaschke]{producto-finito-blaschke}{productos finitos de Blaschke} tales que $\max_{\xi \in \mathbb{T}} \abs{u(\xi) - \frac{B_1(\xi)}{B_2(\xi)}} < \varepsilon$. 
\end{teo}
\begin{dem}
    Por el \excref[lem:fn-continua-toro-imp-vueltas-par-impar]{lem-fn-continua-toro-imp-vueltas-par-impar}, $\exists \eta \in \mathcal{C}(\mathbb{T})$ unimodular $ : \forall \zeta \in \mathbb{T} : u(\zeta) = \eta^2(\zeta)$ o $\forall \zeta \in \mathbb{T} : u(\zeta) = \zeta \cdot \eta^2(\zeta)$. Como $f(\zeta) = \zeta$ es la función frontera de $B \in \mathscr{B}$ dado por $B(z) = z$, en ambos casos basta probar el resultado para $\eta^2$.
    Sin pérdida de generalidad, suponemos que $\varepsilon < 1$.
    
    Por el \exhyperref[teo:stone-weierstrass-toro]{teo-stone-weierstrass-toro}{teorema de Stone--Weierstrass en $\mathbb{T}$},
    $\exists h \colon \mathbb{T} \to \C$ dada por $h(\zeta) = \sum_{k=-N}^N a_k \zeta^k$ tal que $\max_{\xi \in \mathbb{T}} \abs{\eta(\xi) - h(\xi)} < \frac{\varepsilon}{2}$. Entonces, como $\eta$ es unimodular y $\varepsilon < 1$, $h$ no se anula en $\mathbb{T}$. Por tanto, podemos definir $\forall z \in \mathbb{T} : h^* (z) \defeq \bar{h \left(\sfrac{1}{\bar{z}}\right)} \we f(z) \defeq \frac{h(z)}{h^*(z)}$. Entonces, % hay que definir f en todo el plano excepto en los ceros de $h^*$ y el $0$ (el polo de $h$)
    \[
    \begin{aligned}
        \forall \zeta \in \mathbb{T} : \abs{f(\zeta)} = \abs{\frac{h(\zeta)}{h^*(\zeta)}} = \frac{\abs{h(\zeta)}}{\abs{h \left(\sfrac{1}{\bar{\zeta}}\right)}} = \frac{\abs{h(\zeta)}}{\abs{h(\zeta)}} = 1,
    \end{aligned}
    \]
    luego $f$ es unimodular en $\mathbb{T}$. Además, en $\mathbb{T}$, se cumple que
    \[
    \begin{aligned}
        \abs{\eta^2 - f} &= \abs{\eta^2 - \frac{h}{h^*}} = \abs{\frac{\eta}{\eta^*} - \frac{h}{h^*}} = \abs{\frac{\eta h^* - h \eta^* + hh^* - hh^*}{\eta^* h^*}} = \abs{\frac{(\eta - h)h^* + (h^* - \eta^*)h}{\eta^* h^*}} \\
        &\leq \frac{\abs{\eta - h} \cdot \abs{h^*} + \abs{h^* - \eta^*} \cdot \abs{h}}{\abs{\eta^*} \cdot \abs{h^*}} = \frac{\abs{\eta - h} \cdot \abs{h} + \abs{h - \eta} \cdot \abs{h}}{\abs{\eta} \cdot \abs{h}} = 2 \abs{\eta - h} < \varepsilon.
    \end{aligned}
    \]
    Como $f$ es una función meromorfa en $\mathbb{D}$, unimodular y continua en $\mathbb{T}$, por el \excref[cor:carac-cociente-prod-finito-blaschke-algebra-disco-meromorfa]{cor-carac-cociente-prod-finito-blaschke-algebra-disco-meromorfa}, $f$ es un cociente de productos finitos de Blaschke.
\end{dem}

\section{Aproximación por productos finitos de blaschke con ceros simples}

%ORDENAR es innecesario
% \begin{lem}\label{lem:diferencia-productos-leq-suma-diferencias}
%     Sean $\left\{a_k\right\}_{k =1}^n, \left\{b_k\right\}_{k=1}^n \subset \bar{\mathbb{D}}$. Entonces, $\ds \abs{\prod_{k=1}^n a_k - \prod_{k=1}^n b_k} \leq \sum_{k=1}^n \abs{a_k - b_k}$. 
% \end{lem}
% \begin{dem}
%     Lo probaremos por inducción en $n$. Para $n = 1$, el resultado es trivial. Supongamos que el resultado se cumple para $n$. Consideramos $\left\{a_k\right\}_{k =1}^{n+1}, \left\{b_k\right\}_{k=1}^{n+1} \subset \bar{\mathbb{D}}$. Entonces,
%     \[
%     \begin{aligned}
%         \abs{\prod_{k=1}^{n+1} a_k - \prod_{k=1}^{n+1} b_k} &= \abs{\left(\prod_{k=1}^{n+1} a_k - a_{n+1} \prod_{k=1}^{n} b_k + a_{n+1} \prod_{k=1}^{n} b_k - \prod_{k=1}^{n+1} b_k\right)} \\
%         &= \abs{a_{n+1} \left(\prod_{k=1}^{n} a_k - \prod_{k=1}^{n} b_k\right) + \left(\prod_{k=1}^{n} b_k\right) (a_{n+1} - b_{n+1})} \\
%         &\leq \abs{a_{n+1}} \abs{\prod_{k=1}^{n} a_k - \prod_{k=1}^{n} b_k} + \prod_{k=1}^{n} \abs{b_k} \abs{a_{n+1} - b_{n+1}} \\
%         &\upexpl{\leq}{$\!\!\!\!\!\!\!\!\!\!\!\!\!\!\!\!\!\! a_k, b_k \in \bar{\mathbb{D}}$} \abs{\prod_{k=1}^{n} a_k - \prod_{k=1}^{n} b_k} + \abs{a_{n+1} - b_{n+1}} \leq \sum_{k=1}^n \abs{a_k - b_k} + \abs{a_{n+1} - b_{n+1}}.
%     \end{aligned}
%     \]

%     \vspace{-\baselineskip}
%     % \vspace{-\belowdisplayskip}
%     \hfill\qedsymbol
% \end{dem}

\begin{lem}\label{lem:convergencia-ceros-imp-convergencia-prod-finito-blaschke}
    Sea $\left(w_k\right)_{k=1}^n \subset \mathbb{D}$ y sea $B = \prod_{k} \exhyperref[defn:involucion-disco-unidad]{involucion-disco-unidad}{\tau_{w_k}} \in \exhyperref[defn:producto-finito-blaschke]{producto-finito-blaschke}{\mathscr{B}}$. Sea $\left(\underline{z}^m\right)_{m \in \N} \subset \mathbb{D}^n$ tal que
    \[
    \bigl(\forall m \in \N : \underline{z}^m = \left(z_k^m\right)_{k=1}^n \subset \mathbb{D}\bigr) \quad \we \quad \bigl(\forall k \in \N_n : \exhyperref[defn:convergencia]{convergencia}{\lim_{m \to \infty}} z_k^m = w_k\bigr).
    \]
    Definimos $B_m \defeq \prod_{k} \exhyperref[defn:involucion-disco-unidad]{involucion-disco-unidad}{\tau_{z_k^m}} \in \exhyperref[defn:producto-finito-blaschke]{producto-finito-blaschke}{\mathscr{B}}$. Entonces, para todo $K \exhyperref[defn:compacidad]{compacidad}{\Subset} \C$ que no contenga ningún \exhyperref[defn:polo]{polo}{polo} de $B$, se tiene que $B_m \exhyperref[defn:convergencia-uniforme]{convergencia-uniforme}{\xrightrightarrows[m \to \infty]{K}} B$.
    En particular, $B_m \exhyperref[defn:convergencia-uniforme]{convergencia-uniforme}{\xrightrightarrows[m \to \infty]{\bar{\mathbb{D}}}} B$.
\end{lem}
\begin{dem}
    Para cada $k \in \N_n$, calculamos
    \[
    \tau_{w_k}(z) - \tau_{z_k^m}(z) = \frac{\left(w_k - z\right)\left(1 - \bar{z_k^m}z\right) - \left(z_k^m - z\right)\left(1 - \bar{w_k}z\right)}{\left(1 - \bar{w_k}z\right)\left(1 - \bar{z_k^m}z\right)}.
    \]
    Expandiendo el numerador, lo podemos escribir como un polinomio cuadrático en $z$:
    \[
    N_k^m(z) \defeq (w_k - z_k^m) + z\bigl(z_k^m \bar{w_k} - w_k \bar{z_k^m}\bigr) + z^2(\bar{z_k^m} - \bar{w_k}).
    \]
    Como $z_k^m \bar{w_k} - w_k \bar{z_k^m} = (z_k^m - w_k)\bar{w_k} - w_k(\bar{z_k^m} - \bar{w_k})$, se tiene que
    \[
    \abs{N_k^m(z)} \leq \abs{w_k - z_k^m}\left(1 + 2\abs{z}\abs{w_k} + \abs{z}^2\right) \leq \abs{w_k - z_k^m}(1 + \abs{z})^2.
    \]

    Sea $K \Subset \C \setminus \left\{1/\bar{w_k} : k \in \N_n \we w_k \neq 0\right\}$. Entonces, $\delta_k \defeq \min_{z \in K} \abs{1 - \bar{w_k}z} > 0$ están bien definidos. Sea $R \defeq \max_{z \in K} \abs{z}$. Como $z_k^m \to w_k$ y hay un número finito de índices $k$, $\exists m_0 \in \N : \forall m \geq m_0 : \forall k \in \N_n : \abs{z_k^m - w_k} < \frac{\delta_k}{2(R + 1)}$. En particular,
    \[
    \forall z \in K : \abs{1 - \bar{z_k^m}z} \geq \abs{1 - \bar{w_k}z} - \abs{z_k^m - w_k}\abs{z} \geq \delta_k - \frac{\delta_k R}{2(R+1)} \geq \frac{\delta_k}{2},
    \]
    de modo que, para $m \geq m_0$:
    \[
    \sup_{z \in K} \abs{\tau_{w_k}(z) - \tau_{z_k^m}(z)} \leq \frac{(1+R)^2}{\delta_k \cdot \delta_k / 2} \abs{w_k - z_k^m} = \frac{2(1+R)^2}{\delta_k^2} \abs{w_k - z_k^m}.
    \]

    Para pasar del control de cada factor al del producto, procedemos por inducción en $n$. Para $n = 1$, el resultado es inmediato. Supongamos que el resultado se cumple para productos de $n - 1$ factores y definimos $\widetilde{B} = \prod_{k=2}^n \tau_{w_k}$ y $\widetilde{B}_m = \prod_{k=2}^n \tau_{z_k^m}$, de modo que $B = \tau_{w_1} \cdot \widetilde{B}$ y $B_m = \tau_{z_1^m} \cdot \widetilde{B}_m$. Entonces,
    \[
    B(z) - B_m(z) = \tau_{w_1}(z) \bigl(\widetilde{B}(z) - \widetilde{B}_m(z)\bigr) + \bigl(\tau_{w_1}(z) - \tau_{z_1^m}(z)\bigr) \widetilde{B}_m(z).
    \]
    Por hipótesis de inducción, $\widetilde{B}_m \xrightrightarrows[m \to \infty]{K} \widetilde{B}$. Además, como $K$ es compacto y no contiene polos de $B$ ni, para $m$ suficientemente grande, de $B_m$, los factores $\tau_{w_1}$ y $\widetilde{B}_m$ están uniformemente acotados en $K$, i.e. $\exists M < \infty$ tal que $\sup_{z \in K} \abs{\tau_{w_1}(z)} \leq M$ y $\sup_{z \in K} \abs{\widetilde{B}_m(z)} \leq M$ para $m$ suficientemente grande. Por tanto,
    \[
    \sup_{z \in K} \abs{B(z) - B_m(z)} \leq M \sup_{z \in K} \abs{\widetilde{B}(z) - \widetilde{B}_m(z)} + M \sup_{z \in K} \abs{\tau_{w_1}(z) - \tau_{z_1^m}(z)} \xrightarrow[m \to \infty]{} 0.
    \]

    \vspace{-\baselineskip}
    % \vspace{-\belowdisplayskip}
    \hfill\qedsymbol
\end{dem}

\begin{teo}\label{teo:aproximacion-prod-finito-blaschke-ceros-simples}
    Sea $B \in \exhyperref[defn:producto-finito-blaschke]{producto-finito-blaschke}{\mathscr{B}_n}$. Entonces, $\exists \left(B_m\right)_{m \in \N} \subset \exhyperref[defn:producto-finito-blaschke]{producto-finito-blaschke}{\mathscr{B}_n}$  tal que:
    \begin{enumerate}[ref=(\arabic*)]
        \item\label{teo:aproximacion-prod-finito-blaschke-ceros-simples:ceros-simples} $\forall m \in \N : B_m$ tiene \exhyperref[defn:orden-cero-fn-holomorfa]{orden-cero-fn-holomorfa}{ceros simples}.
        \item\label{teo:aproximacion-prod-finito-blaschke-ceros-simples:origen} $\forall m \in \N : B_m(0) \neq 0 \we B_m'(0) \neq 0$.
        \item\label{teo:aproximacion-prod-finito-blaschke-ceros-simples:convergencia} $\forall K \exhyperref[defn:compacidad]{compacidad}{\Subset} \C : K$ no contiene ningún \exhyperref[defn:polo]{polo}{polo} de $B : B_m \exhyperref[defn:convergencia-uniforme]{convergencia-uniforme}{\xrightrightarrows[m \to \infty]{K}} B$.
    \end{enumerate}
\end{teo}
\begin{dem}
    Por definición, $\exists \left\{z_k\right\}_{k=1}^r \subset \mathbb{D}, \left\{j_k\right\}_{k=1}^r \subset \N, j_0 \in \N \cup \{0\}, \beta \in [0, 2\pi) : $ tales que $B$ viene dada por
    \[
    B(z) = e^{i\beta} z^{j_0} \prod_{k=1}^r \left(\frac{z_k - z}{1 - \bar{z_k}z}\right)^{j_k} \quad \tex{donde} \quad \sum_{k=0}^r j_k = n
    \]
    y $\left\{z_k\right\}_{k=1}^r$ son los distintos ceros de $B$ en $\mathbb{D} \setminus \{0\}$. Definiendo $z_0 = 0$, consideramos
    \[
    \varepsilon_0 \defeq \min \left\{ \frac{1}{2} \abs{z_k - z_\ell} : 0 \leq k, \ell \leq r \we k \neq \ell \right\} > 0.
    \]
    Entonces, existe un $N$ tal que $\frac{1}{N} < \varepsilon_0$. Fijamos $m \in \N$. Definimos los discos
    \[
    \forall k \in \N_r \cup \{0\} : \Gamma_{k, m} \defeq \left\{z \in \mathbb{C} : \abs{z - z_k} < \frac{1}{m+N} \right\}
    \]
    que son disjuntos dos a dos y no contienen al origen.
    
    Para cada $m \in \N$ y $k \in \N_r \cup \{0\}$, tomamos $\forall \ell \in \N_{j_k} : z_{m, k, \ell} = z_k + \frac{1}{(m+N)} \cdot \omega_{j_k}^\ell$ donde $\omega_{j_k}$ es la raíz $j_k$-ésima de la unidad. Entonces, $\forall k \in \N_r \cup \{0\} : \left\{z_{m, k, \ell} \right\}_{\ell=1}^{j_k} \subset \Gamma_{k, m}$ y los $z_{k, m, \ell}$ son distintos entre sí. Para $m \in \N$, consideramos $B_m \in \mathscr{B}_n$ mónico dado por
    \[
    B_m(z) = e^{i\beta} \prod_{k=0}^r \prod_{\ell=1}^{j_k} \frac{z_{m, k, \ell} - z}{1 - \bar{z_{m, k, \ell}} z}.
    \]
    Las condiciones \ref{teo:aproximacion-prod-finito-blaschke-ceros-simples:ceros-simples} y la primera parte de \ref{teo:aproximacion-prod-finito-blaschke-ceros-simples:origen} se cumplen por construcción. La condición \ref{teo:aproximacion-prod-finito-blaschke-ceros-simples:convergencia} se sigue del \excref[lem:convergencia-ceros-imp-convergencia-prod-finito-blaschke]{lem-convergencia-ceros-imp-convergencia-prod-finito-blaschke}.
    Para comprobar si $B_m'(0) \neq 0$, reindexamos los $\left\{z_{k, m, \ell}\right\}_{k=1, \ell=1}^{r, j_k}$ como $\left\{w_p\right\}_{p=1}^{n}$ para obtener: $B_m = e^{i\beta} \prod_{p=1}^{n} \tau_{w_p}$. Entonces, por el \excref[lem:derivada-logaritmica-prod-finito-blaschke]{lem-derivada-logaritmica-prod-finito-blaschke}, se tiene que 
    \[
    \frac{B_m'(0)}{B_m(0)} = \sum_{p=1}^n \frac{\abs{w_p}^2 - 1}{w_p} = \sum_{p=1}^n \left(\bar{w_p} - \frac{1}{w_p}\right) = \left(\bar{w_n} - \frac{1}{w_n}\right) + \sum_{p=1}^{n-1} \left(\bar{w_p} - \frac{1}{w_p}\right).
    \]
    En caso de que $B_m'(0) = 0$ y el sumario anterior se anule, tomamos $w_{n} = z_{m, r, j_r} = z_r + \frac{1}{m+N} \cdot \gamma$ con $\gamma \in \mathbb{T}$ apropiadamente elegido en lugar de la expresión original, de modo que $B_m'(0) \neq 0$.%REVISAR y COMPLETAR
\end{dem}

\section{Generalización del Teorema de Rouché y su inverso}

\begin{defn}\label{defn:fn-holomorfa-dentro-camino}
    Sea $\gamma \colon [0, 1] \to \C$ un camino cerrado, $f$ es holomorfa dentro de $\gamma$
    \[
    \iff \exists \Omega \subset \C \tex{ abierto y simplemente conexo} : \gamma([0, 1]) \subset \Omega \we f \in \mathcal{H}(\Omega).
    \]
\end{defn}

\begin{defn}[Polo]\label{defn:polo-fn-meromorfa}
    Sea $\Omega \subset \C$ abierto, $z_0 \in \Omega$ y $f \in \mathcal{H}(\Omega \setminus \{z_0\})$,
    \[
    z_0 \tex{ es un polo de } f \iff \lim_{z \to z_0} \abs{f(z)} = \infty.
    \]
\end{defn}

\begin{defn}[Orden de un polo]\label{defn:orden-polo-fn-meromorfa}
    Sea $z_0$ un polo de $f$, $\operatorname{ord}(f, z_0)$ es el orden de $z_0$
    \[
    \iff \operatorname{ord}(f, z_0) = \min \left\{ n \in \N : \lim_{z \to z_0} \abs{(z - z_0)^n f(z)} < \infty \right\}.
    \]
\end{defn}

\begin{defn}\label{defn:numero-ceros-fn-holomorfa-camino}
    Sea $\Gamma \colon [0, 1] \to \C$ un camino cerrado y $f$ holomorfa dentro de $\Gamma$, $Z_f(\Gamma)$ es el número de ceros de $f$ dentro de $\Gamma$ contado con multiplicidad
    \[
    \iff Z_f(\Gamma) = \sum_{\substack{z \in \operatorname{int}(\Gamma) \\ f(z) = 0}} \operatorname{ord}(f, z).
    \]
\end{defn}

\begin{defn}\label{defn:numero-polos-fn-meromorfa-camino}
    Sea $\Gamma \colon [0, 1] \to \C$ un camino cerrado y $f$ meromorfa dentro de $\Gamma$, $P_f(\Gamma)$ es el número de polos de $f$ dentro de $\Gamma$ contado con multiplicidad
    \[
    \iff P_f(\Gamma) = \sum_{\substack{z \in \operatorname{int}(\Gamma) \\ z \tex{ es un polo de } f}} \operatorname{ord}(f, z).
    \]
\end{defn}

\begin{teo}[Principio del argumento]\label{teo:principio-argumento}
    Sea $\Omega \subset \C$ un dominio, $\Gamma \colon [0, 1] \to \Omega$ un camino simple y cerrado orientado positivamente. Sea $f$ meromorfa en $\Omega$ y sin ceros ni polos en $\Gamma$. Entonces,
    \[
    Z_f(\Gamma) - P_f(\Gamma) = \frac{1}{2\pi i} \int_{\Gamma} \frac{f'(z)}{f(z)} \odif{z}.
    \]
\end{teo}%DEMOSTRACIÓN

\begin{teo}[de Rouché]\label{teo:rouche}
    Sea $\Omega \subset \C$ un abierto, $f, g \in \mathcal{H}(\Omega)$ y sea $\gamma$ una \exhyperref[defn:curva-jordan]{Curva-jordan}{curva de Jordan} rectificable contenida en $\Omega$ tal que su interior $\operatorname{int}(\gamma) \subset \Omega$. Si
    \[
    \forall z \in \gamma : \abs{f(z) - g(z)} < \abs{f(z)},
    \]
    entonces $f$ y $g$ tienen el mismo número de ceros en $\operatorname{int}(\gamma)$ (contando multiplicidades).
\end{teo}
\begin{dem}
    Supongamos por contradicción que $\exists z_0 \in \gamma$ tal que $f(z_0) = 0$. Entonces, $\abs{g(z_0)} = \abs{f(z_0) - g(z_0)} < \abs{f(z_0)} = 0$, lo cual es imposible.
    
    Además, por la desigualdad triangular inversa, $\forall z \in \gamma : \abs{g(z)} \geq \abs{f(z)} - \abs{f(z) - g(z)} > 0$, luego ni $f$ ni $g$ tienen ceros en $\gamma$.

    Consideramos $h(z) = 1 - \frac{g(z)}{f(z)}$. Entonces, $h \in \mathcal{H} \left(\Omega \setminus \left\{z \in \Omega : f(z) = 0\right\}\right)$ y se cumple que
    \[
    \forall z \in \gamma : \abs{h(z)} = \abs{1 - \frac{g(z)}{f(z)}} = \abs{\frac{f(z) - g(z)}{f(z)}} = \frac{\abs{f(z) - g(z)}}{\abs{f(z)}} < 1.
    \]
    Por tanto, $h(\gamma) \subset \mathbb{D}$. En particular, la curva cerrada $h \circ \gamma$ está contenida en $\mathbb{D}$ y no pasa por el origen (pues $h$ no se anula en $\gamma$, ya que $h(z) = 0 \iff g(z) = f(z)$, lo cual implicaría $\abs{f(z) - g(z)} = 0 < \abs{f(z)}$ $\contr$).%REVISAR demostración a partir de aquí, consultar fuentes.
    
    Como $h \circ \gamma$ es una curva cerrada en $\mathbb{D} \setminus \{0\}$ y el origen no está en la imagen, podemos considerar una homotopía continua que contrae $h \circ \gamma$ al punto $h(\gamma(0)) \in \mathbb{D}$ sin pasar por el origen. Por tanto, el índice de $h \circ \gamma$ respecto al origen es cero, es decir,
    \[
    \frac{1}{2\pi i} \int_\gamma \frac{h'(z)}{h(z)} \odif{z} = 0.
    \]

    Ahora bien, como $h(z) = 1 - \sfrac{g(z)}{f(z)} = \frac{f(z) - g(z)}{f(z)}$, tenemos que
    \[
    \frac{h'(z)}{h(z)} = \frac{\odv{}{z} \left(\frac{f(z) - g(z)}{f(z)}\right)}{\frac{f(z) - g(z)}{f(z)}} = \frac{f(z)(f'(z) - g'(z)) - (f(z) - g(z))f'(z)}{f(z) - g(z)} = \frac{f'(z)}{f(z)} - \frac{g'(z)}{g(z)}.
    \]

    Por tanto,
    \[
    0 = \frac{1}{2\pi i} \int_\gamma \frac{h'(z)}{h(z)} \odif{z} = \frac{1}{2\pi i} \int_\gamma \frac{f'(z)}{f(z)} \odif{z} - \frac{1}{2\pi i} \int_\gamma \frac{g'(z)}{g(z)} \odif{z}.
    \]

    Por el principio del argumento para funciones holomorfas, tenemos que $Z_f(\gamma) = Z_g(\gamma)$. Es decir, $f$ y $g$ tienen el mismo número de ceros dentro de $\gamma$ contando multiplicidades.
\end{dem}

\begin{teo}[Glicksberg]\label{teo:glicksberg}
    Sea $\Gamma \colon [0, 1] \to \C$ un camino cerrado simple y $f, g$ holomorfas dentro de $\Gamma$ tales que $\forall z \in \Gamma : \abs{f(z) - g(z)} < \abs{f(z)} + \abs{g(z)}$. Entonces, $f$ y $g$ tienen el mismo número de ceros dentro de $\Gamma$ contando multiplicidades.
\end{teo}
\begin{dem}
    Supongamos por contradicción que $\exists z_0 \in \Gamma$ tal que $f(z_0) = 0$. Entonces, $\abs{g(z_0)} = \abs{f(z_0) - g(z_0)} < \abs{f(z_0)} + \abs{g(z_0)} = \abs{g(z_0)}$, lo cual es imposible. Por tanto, $f$ no se anula en $\Gamma$. De forma análoga, $g$ no se anula en $\Gamma$. Entonces, sabemos que
    \[
    \forall z \in \Gamma : \abs{\frac{f(z)}{g(z)} - 1} < \abs{\frac{f(z)}{g(z)}} + 1.
    \]
    Esto implica que $\forall z \in \Gamma : \sfrac{f(z)}{g(z)} > 0$ (si $\sfrac{f(z)}{g(z)} \leq 0$, entonces $1 \leq \abs{\sfrac{f(z)}{g(z)} - 1} < \abs{\sfrac{f(z)}{g(z)}} + 1 \leq 1$, lo cual es imposible). %REVISAR esto está mal mirar la foto del 03/03/2026
    Por tanto, $\sfrac{f}{g} \left(\Gamma\right) \subset \C \setminus (-\infty, 0]$. Entonces, tomando la rama principal del logaritmo, $\log \left(\sfrac{f}{g}\right) \in \mathcal{H} \left(\operatorname{int}(\Gamma)\right)$ y, al derivar, se tiene que
    \[
    \odv{}{z} \log \left(\frac{f(z)}{g(z)}\right) = \frac{\odv{}{z} \left(\frac{f(z)}{g(z)}\right)}{\frac{f(z)}{g(z)}} = \frac{f'(z)g(z) - f(z)g'(z)}{f(z)g(z)} = \frac{f'(z)}{f(z)} - \frac{g'(z)}{g(z)}.
    \]

    Por el principio del argumento para funciones holomorfas, como $Z_{\sfrac{f}{g}} (\Gamma) = 0$, tenemos que
    \[
    0 = \frac{1}{2\pi i} \int_\Gamma \odv{}{z} \log \left(\frac{f(z)}{g(z)}\right) \odif{z} = \frac{1}{2\pi i} \int_\Gamma \frac{f'(z)}{f(z)} \odif{z} - \frac{1}{2\pi i} \int_\Gamma \frac{g'(z)}{g(z)} \odif{z}.
    \]
    De nuevo por el principio del argumento, $Z_f(\Gamma) = Z_g(\Gamma)$, es decir, $f$ y $g$ tienen el mismo número de ceros dentro de $\Gamma$ contando multiplicidades.
\end{dem}

\begin{obs}
    El teorema de Glicksberg es una generalización del teorema de Rouché, ya que la hipótesis de Rouché implica la de Glicksberg: si $\abs{f(z) - g(z)} < \abs{f(z)}$, entonces
    \[
    \abs{f(z) - g(z)} < \abs{f(z)} \leq \abs{f(z)} + \abs{g(z)}.
    \]

    Sin embargo, la recíproca es falsa en general. Por ejemplo, si tomamos el camino simple cerrado $\Gamma = \mathbb{T}$ y las funciones enteras $f(z) = 2z$ y $g(z) = iz$, entonces
    \[
    \forall z \in \mathbb{T} : \abs{f(z) - g(z)} = \abs{(2-i)z} = \sqrt{5} \;\;\we\;\; \abs{f(z)} = 2 \;\;\we\;\; \abs{f(z)} + \abs{g(z)} = 3
    \]
    Como $\sqrt{5} < 3$, la hipótesis de Glicksberg se cumple. Sin embargo, como $\sqrt{5} > 2$, la hipótesis de Rouché falla. En todo caso, la conclusión es válida: $f$ y $g$ tienen ambas un único cero en $\mathbb{D}$ (en $z = 0$).
\end{obs}

\begin{obs}
    Sean $f, g \in \mathcal{H} \left(D(0, R)\right)$ con $R> 1$ y sin ceros en $\mathbb{T}$. Sean $B_1, B_2$ productos finitos de Blaschke de grado $n$ tales que $\abs{B_1 f + B_2 g} < \abs{f} + \abs{g}$ en $\mathbb{T}$. Entonces, $\abs{B_1 f - B_2 g} < \abs{B_1 f} + \abs{B_2 g}$ en $\mathbb{T}$, luego por el teorema de Glicksberg, %ORDENAR referenciar
    $Z_{B_1 f}(\mathbb{T}) = Z_{B_2 g}(\mathbb{T})$. Es decir, $n + Z_f(\mathbb{T}) = n + Z_g(\mathbb{T})$, luego $Z_f(\mathbb{T}) = Z_g(\mathbb{T})$.

    Como muestra el siguiente teorema, el recíproco de este resultado también es cierto.
\end{obs}

\begin{teo}[Challener--Rubel]\label{teo:challener-rubel}
    Sean $f, g \in \mathcal{H}\left(D(0,R)\right)$ con $R > 1$ y sin ceros en $\mathbb{T}$. Entonces, $Z_f(\mathbb{T}) = Z_g(\mathbb{T})$ si y solo si $\exists B_1, B_2$ productos finitos de Blaschke del mismo grado tal que $\abs{B_1 f + B_2 g} < \abs{f} + \abs{g}$ en $\mathbb{T}$.
\end{teo}
\begin{dem}
    La implicación $\boxed{\Leftarrow}$ es la observación anterior. %ORDENAR referenciar

    Para la implicación $\boxed{\Rightarrow}$, como $f$ y $g$ son funciones continuas sin ceros en $\mathbb{T}$, 
    \[
    m = \min_{\xi \in \mathbb{T}} \min \left\{\abs{f(\xi)}, \abs{g(\xi)}\right\} > 0  \quad \we \quad M = \max_{\xi \in \mathbb{T}} \max \left\{\abs{f(\xi)}, \abs{g(\xi)}\right\} < \infty.
    \]
    Definimos por tanto $u \colon \mathbb{T} \to \mathbb{T}$ dada por $u(\xi) = \frac{g/f}{\abs{g/f}}(\xi)$. Entonces, $u$ es continua y, por el teorema de Helson--Sarason aplicado a $-u$, %ORDENAR referenciar
    $\exists B_1, B_2$ productos finitos de Blaschke tales que 
    \[
    \max_{\xi \in \mathbb{T}} \abs{u(\xi) + \frac{B_1(\xi)}{B_2(\xi)}} < \frac{m}{M}.
    \]

    Entonces, en $\mathbb{T}$, se cumple que
    \[
    \begin{aligned}
        \abs{B_1 f + B_2 g} &= \abs{f} \cancelto{1}{\abs{B_2}}\quad \abs{\frac{B_1}{B_2} + \frac{g}{f}} = \abs{f} \cdot \abs{u \abs{\frac{g}{f}} + \frac{B_1}{B_2}} \\
        &= \abs{f} \cdot \abs{u \left(\abs{\frac{g}{f}} -1\right) + \left(u + \frac{B_1}{B_2}\right)} \leq \abs{f} \cdot \cancelto{1}{\abs{u}}\; \cdot \abs{\frac{\abs{g}}{\abs{f}} - 1} + \abs{f} \cdot \abs{u + \frac{B_1}{B_2}} \\
        &<  \abs{\abs{g} - \abs{f}} + M \cdot \frac{m}{M} = \abs{\abs{f} - \abs{g}} + m \leq \abs{f} + \abs{g}. 
    \end{aligned}
    \]
    Donde la última desigualdad se sigue de la identidad $\abs{f} + \abs{g} - \abs{\abs{f} - \abs{g}} = 2 \min \left\{\abs{f}, \abs{g}\right\}$ y de la definición de $m$. %ORDENAR: a lo mejor esto no hace falta explicarlo

    Como $\abs{B_1 f} + \abs{B_2 g} = \abs{f} + \abs{g}$ en $\mathbb{T}$, se sigue que $\abs{B_1 f + B_2 g} < \abs{B_1 f} + \abs{B_2 g}$ en $\mathbb{T}$, luego por el teorema de Glicksberg, %ORDENAR referenciar
    $Z_{B_1 f}(\mathbb{T}) = Z_{B_2 g}(\mathbb{T})$ y como $Z_{f} (\mathbb{T}) = Z_{g}(\mathbb{T})$, se tiene que $Z_{B_1}(\mathbb{T}) = Z_{B_2}(\mathbb{T})$, es decir, $B_1$ y $B_2$ tienen el mismo grado.
\end{dem}


\newpage

\lipsum[1-15]
